\documentclass[11pt]{article}

\input{../shared/preamble.tex}

\usepackage{amsmath}
\usepackage{algorithm}
\usepackage{algpseudocode}
\usepackage{cleveref}

\title{Thurstone Portfolios:\\
       Long-Only Allocation by Inverting Winning Probabilities}
\author{Peter Cotton}
\date{\today}

\newcommand{\E}{\mathbb{E}}
\newcommand{\Prob}{\mathbb{P}}
\newcommand{\R}{\mathbb{R}}
\DeclareMathOperator*{\argminop}{arg\,min}
\DeclareMathOperator*{\argmaxop}{arg\,max}

\newtheorem{theorem}{Theorem}
\newtheorem{proposition}{Proposition}

\begin{document}
\maketitle

\begin{abstract}
We develop \emph{Thurstone portfolios}, a long-only allocation scheme in which
portfolio weights are the winning probabilities of a Thurstonian race among
assets. The construction uses two correlation matrices: we calibrate latent
\emph{abilities} so that the race under a reference correlation reproduces a
chosen benchmark (such as capitalization weights), then re-run the race under the
actual estimated correlation and read off the winners as the tilted weights. The
portfolio is thus a \emph{controlled perturbation} of the benchmark---it
reproduces the benchmark exactly when the two correlations agree, and departs
from it only in response to their difference. The weights lie on the simplex by
construction (no covariance inversion, no constraints), solve a
correlation-aware regularized allocation problem, and---unlike capitalization
weighting, which is a \emph{Lucian} (proportional-choice) scheme---are immune to
the duplication (independence-of-irrelevant-alternatives) paradox: perfectly
correlated holdings share a single position's weight. We prove the weights vary
smoothly (locally Lipschitz) with the correlation, which bounds turnover and
supports cheap online rebalancing through a common-seed transport that scales to
thousands of assets. We position the method against the Capital Asset Pricing
Model (CAPM), whose implied market-capitalization weighting we argue is a Lucian
special case a Thurstonian model need not reproduce, and outline a backtesting
protocol with two empirical studies (REITs and Russell-index equities tilted by
language-model ``AI-benefit'' scores). The method is implemented in the
\texttt{allocation} package, built on \texttt{thurstone}.
\end{abstract}

\section{Introduction}
\label{sec:intro}

Modern passive investing rests on an equilibrium argument: under the CAPM
\parencite{sharpe1964}, the market-capitalization-weighted portfolio is mean--variance
efficient, so a rational investor should simply hold the index. There is a less
remarked structural assumption underneath. Capitalization weighting treats the
relative ``attractiveness'' of two assets as a fixed ratio that is preserved
when other assets are added to or removed from the universe. This is precisely
Luce's choice axiom \parencite{luce1959}, the independence-of-irrelevant-alternatives
property that underlies the Bradley--Terry \parencite{bradleyterry1952} and
Plackett--Luce \parencite{plackett1975} models of choice and ranking. We will call a
weighting scheme with this property \emph{Lucian}.

The Lucian assumption is convenient but neither logically necessary nor
empirically robust. Investors do not observe the true covariance of returns;
they estimate it, and good practice is to \emph{shrink} those estimates toward a
structured target \parencite{ledoitwolf2004}. Shrinkage breaks the
wisdom-of-crowds aggregation that would make the capitalization-weighted
portfolio efficient, and tracking the index then carries a cost in
risk-adjusted return unless the index is itself optimal \parencite{roll1992}.
This invites a question: if the market need not reveal itself as exactly Lucian,
what is a natural non-Lucian alternative?

Thurstone's law of comparative judgment \parencite{thurstone1927} supplies one.
In a Thurstonian model each asset $i$ carries a latent \emph{ability} $a_i$ and a
noisy \emph{performance} $X_i = a_i + \varepsilon_i$; the asset with the extreme
performance ``wins.'' Unlike the Lucian model, the probability that $i$ beats $j$
depends on the rest of the field and---crucially for finance---on the
\emph{correlation} of the performances. The contribution of this paper is to use
the Thurstonian winning probabilities themselves as portfolio weights, and to
show that the inverse problem (recovering abilities from a target allocation, and
re-deriving weights once correlation is injected) admits a fast, stable solution.

\paragraph{Contributions.}
\begin{enumerate}
  \item \textbf{Construction.} We define the \emph{ability tilt}: calibrate
        latent abilities so that a Thurstonian race under a reference correlation
        reproduces a chosen benchmark, then re-run the race under the estimated
        correlation. The weights are long-only by construction, and the portfolio
        is a controlled perturbation of the benchmark---identical to it when the
        two correlations agree (\Cref{sec:method}, \Cref{prop:fixedpoint}).
  \item \textbf{Non-Lucian behaviour.} We prove the tilt is
        redundancy-coherent---perfectly correlated holdings share a single
        position's weight (\Cref{prop:redundancy})---and that a correlated
        subgroup's combined weight is monotone decreasing in its internal
        correlation (\Cref{prop:monotone}, via Slepian's inequality), with no
        explicit risk term.
  \item \textbf{Implied objective.} We identify the objective the weights
        optimize: $w = \nabla G_C(\theta)$ for the Gaussian expected-maximum
        $G_C$, equivalently a correlation-aware regularized program on the simplex
        (\Cref{thm:objective}); capitalization weighting is the
        correlation-blind (Gumbel/softmax) special case.
  \item \textbf{Smoothness and turnover.} We prove the weights are $C^\infty$ and
        locally Lipschitz in the correlation (\Cref{thm:smooth}, via Price's
        theorem), and give a common-seed transport whose turnover is bounded by
        the change in the correlation rather than by sampling noise, with a
        factor-form route to large universes (\Cref{sec:smoothness},
        \Cref{sec:computation}).
  \item \textbf{CAPM as a Lucian special case.} We recast Roll's critique in
        choice-theoretic terms: capitalization weighting is \emph{Lucian} and
        therefore inconsistent with efficiency on any restricted universe
        (\Cref{prop:restriction}). We state the \emph{Indispensable Markets
        Hypothesis} and outline a backtesting protocol with two empirical studies
        (\Cref{sec:capm}, \Cref{sec:experiments}).
\end{enumerate}

\section{The Thurstonian Model of a Field}
\label{sec:model}

Consider a universe of $n$ assets. Each asset has a latent ability $a_i\in\R$ and
a performance
\begin{equation}
  X_i = a_i + \varepsilon_i, \qquad (\varepsilon_1,\dots,\varepsilon_n) \sim
  \mathcal{N}(0, \Sigma_\varepsilon),
  \label{eq:performance}
\end{equation}
where we adopt the convention that the \emph{minimum} performance wins, so that a
smaller ability denotes a stronger competitor. The \emph{state price} (winning
probability) of asset $i$ is
\begin{equation}
  p_i \;=\; \Prob\!\left( X_i = \min_{k} X_k \right).
  \label{eq:stateprice}
\end{equation}
When the performances are independent ($\Sigma_\varepsilon$ diagonal) the
$p_i$ are determined by the abilities alone through the distribution of the first
order statistic, and $p$ is a smooth, monotone function of $a$
\parencite{cotton2021}. Two maps are then of interest:
\begin{align}
  \text{forward:}\quad & a \;\longmapsto\; p, \label{eq:forward}\\
  \text{inverse:}\quad & p \;\longmapsto\; a. \label{eq:inverse}
\end{align}
The inverse \eqref{eq:inverse}---the \emph{Fast Ability Transform}---recovers
abilities (up to an additive constant) from a target set of winning
probabilities, and is the engine of the construction below.

\section{Thurstone Portfolios via the Ability Tilt}
\label{sec:method}

We treat portfolio \emph{weights} as winning probabilities and \emph{assets} as
competitors. The construction uses \emph{two} correlation matrices. We first
\emph{calibrate} latent abilities so that the race under a reference correlation
$C_{\mathrm{calib}}$ reproduces a target allocation $w^{\mathrm{target}}$ (a
benchmark such as capitalization weights); we then re-run the race under a second
correlation $C_{\mathrm{tilt}}$---the actual estimated correlation of
returns---and read off the winning probabilities as the tilted weights. The
abilities pin the portfolio to the benchmark at the reference correlation; the
tilt is the portfolio's response to $C_{\mathrm{tilt}} - C_{\mathrm{calib}}$
(\Cref{prop:fixedpoint}).

\begin{algorithm}[h]
\caption{Ability tilt (Thurstone portfolio)}
\label{alg:tilt}
\begin{algorithmic}[1]
\Require target weights $w^{\mathrm{target}}\in\Delta$; reference correlation
         $C_{\mathrm{calib}}$; tilt correlation $C_{\mathrm{tilt}}$.
\State $a \gets \textsc{Calibrate}(w^{\mathrm{target}}, C_{\mathrm{calib}})$
       \Comment{abilities s.t.\ the race under $C_{\mathrm{calib}}$ yields $w^{\mathrm{target}}$}
\State Draw $X^{(1)},\dots,X^{(M)} \sim \mathcal{N}(a, C_{\mathrm{tilt}})$
\State $w_i \gets \dfrac{1}{M}\sum_{m=1}^{M}
       \mathbf{1}\!\left[\, i = \argminop_k X^{(m)}_k \,\right]$
       \Comment{win frequency}
\State \Return $w$
\end{algorithmic}
\end{algorithm}

Two choices of the reference correlation give the practical flavours:
\begin{itemize}
  \item \textbf{Diagonal} ($C_{\mathrm{calib}} = I$). Calibration is the exact
        order-statistic inverse of \Cref{sec:model} (the Fast Ability
        Transform), and the tilt then injects the estimated correlation. The
        target may be any benchmark---the variance-only diagonal portfolio,
        capitalization weights, or equal weight.
  \item \textbf{One-factor market}
        ($C_{\mathrm{calib}} = \beta\beta^\top + \mathrm{diag}(1-\beta^2)$). The
        reference is a single market factor with loadings $\beta$; calibration
        reproduces the benchmark \emph{given} that market correlation, so the
        tilt's deviation is driven by the residual (non-market) correlation in
        $C_{\mathrm{tilt}}$.
\end{itemize}

\noindent A scalar $\phi \in [0,1]$ may be used to set
$C_{\mathrm{tilt}} = (1-\phi)\,C_{\mathrm{calib}} + \phi\,\widehat{C}$, dialling
the estimated correlation $\widehat{C}$ in continuously: $\phi = 0$ reproduces the
benchmark (\Cref{prop:fixedpoint}) and $\phi = 1$ uses the full estimate. The
construction never inverts a covariance---it needs only a valid correlation for
sampling, repaired if necessary to the nearest correlation matrix, so it tolerates
ill-conditioned or singular estimates. Its feasibility, fixed-point, redundancy,
monotonicity, and objective properties are established in \Cref{sec:properties}.

\section{Computation}
\label{sec:computation}

\paragraph{Calibration (deterministic).}
The calibration step---the inverse map from target weights to abilities---is
computed deterministically with the lattice machinery of the \texttt{thurstone}
package, which builds the distribution of the first order statistic by convolving
competitor densities and tracking conditional multiplicity for ties
\parencite{cotton2021}. For a diagonal reference this is the exact
order-statistic inverse \eqref{eq:inverse}, with no Monte-Carlo error and cost
linear in $n$. For the one-factor reference the assets are conditionally
independent given the market factor, so the same lattice race is evaluated at a
few Gauss--Hermite nodes and integrated over the factor; the inverse is then a
damped fixed-point on this forward map. This is the only place quadrature enters.

\paragraph{Tilt (Monte-Carlo).}
The correlated race under the full $C_{\mathrm{tilt}}$ has no tractable analytic
form, so it is evaluated by sampling: performances are drawn from
$\mathcal{N}(a, C_{\mathrm{tilt}})$---with a randomized Sobol' quasi-Monte-Carlo
sequence when available---and the weights are the argmin frequencies. This is the
construction of the \texttt{allocation} package, which supersedes the earlier
\texttt{winningport} implementation of the same tilt.

\paragraph{Online updates and scale.}
At rebalancing the correlation estimate moves slowly, and \Cref{sec:smoothness}
shows the weights move smoothly with it. We exploit this by holding the
standard-normal seeds fixed and \emph{transporting} the path ensemble to the new
correlation (recolouring by the updated square root), so turnover tracks genuine
correlation change rather than sampling noise. For large universes the dense
correlation is never formed: a low-rank factor model $C = BB^\top + D$ keeps both
the transport and the argmin streaming and linear in $n$, scaling to thousands of
assets. The covariance estimate itself is supplied by any external online
estimator, so the method interoperates with standard covariance-forecasting
tooling.

\section{Theoretical Properties}
\label{sec:properties}

We collect the properties that make the ability tilt a principled allocator
rather than a heuristic. Throughout, the abilities $a$ are fixed and we view the
weights as the winning-probability map $w(C)$.

\begin{proposition}[Feasibility]
\label{prop:feasible}
For every positive-definite $C$ and every $a$, the weights lie in the
probability simplex: $w_i \ge 0$ and $\sum_i w_i = 1$.
\end{proposition}
\begin{proof}
With fair tie-splitting exactly one winner is assigned in each realization, so
the $w_i$ are the probabilities of a partition of the sample space.
\end{proof}
\noindent Long-only and fully-invested portfolios thus come for free---no
constraints, penalties, or projection onto the simplex.

\begin{proposition}[Index reproduction]
\label{prop:fixedpoint}
If the tilt correlation equals the calibration correlation,
$C_{\mathrm{tilt}} = C_{\mathrm{calib}}$, then $w = w^{\mathrm{target}}$.
\end{proposition}
\begin{proof}
The abilities are calibrated so that the race under $C_{\mathrm{calib}}$
reproduces the target; evaluating at $C_{\mathrm{tilt}} = C_{\mathrm{calib}}$
returns it.
\end{proof}
\noindent The tilt is therefore a \emph{controlled perturbation} of the
benchmark: the deviation is driven by $C_{\mathrm{tilt}} - C_{\mathrm{calib}}$
and bounded by the Lipschitz estimate of \Cref{thm:smooth}.

\begin{proposition}[Redundancy consistency]
\label{prop:redundancy}
Replace asset $k$ by two perfect duplicates $k', k''$ (equal ability, unit
variance, mutual correlation $1$, and the same correlation as $k$ with every
other asset). Then, with fair tie-splitting,
\begin{equation*}
  w_{k'} + w_{k''} = w_k^{\setminus}, \qquad
  w_{k'} = w_{k''} = \tfrac12 w_k^{\setminus}, \qquad
  w_j = w_j^{\setminus}\ \ (j \neq k', k''),
\end{equation*}
where $w^{\setminus}$ is the allocation in the original field containing the
single asset $k$.
\end{proposition}
\begin{proof}
Perfectly correlated unit-variance Gaussians with equal mean are almost surely
equal, $X_{k'} = X_{k''}$ a.s., so the field minimum---and hence every other
asset's winning event---is unchanged, while the pair ties exactly when $k$ would
have won. Fair splitting halves that shared probability.
\end{proof}

This is the rigorous form of the non-Lucian advantage. At the independent
extreme a duplicate is a \emph{fresh} draw, so the pair's combined weight
strictly exceeds $w_k^{\setminus}$ and draws share from every other asset---the
independence-of-irrelevant-alternatives (``red bus / blue bus'') paradox that
capitalization weighting inherits from Luce's axiom (\Cref{sec:capm}). As the
duplicate's correlation rises from $0$ to $1$, \Cref{thm:smooth} guarantees the
pair's share moves \emph{continuously} from the Lucian value to the
redundancy-coherent $w_k^{\setminus}$. \Cref{tab:redundancy} shows this for
three equal-ability assets.

\begin{table}[h]
  \centering
  \caption{Redundancy consistency and diversification monotonicity. Three
  equal-ability assets; assets $1,2$ are duplicates with mutual correlation
  $\rho_{12}$, asset $0$ independent of them ($M = 4\times10^5$). The duplicated
  pair's share falls monotonically from the Lucian $2/3$ to the
  redundancy-coherent $1/2$.}
  \label{tab:redundancy}
  \begin{tabular}{ccc}
    \toprule
    $\rho_{12}$ & $w_0$ & $w_1 + w_2$ \\
    \midrule
    $0.00$ & $0.333$ & $0.667$ \\
    $0.50$ & $0.384$ & $0.616$ \\
    $0.90$ & $0.449$ & $0.551$ \\
    $0.99$ & $0.484$ & $0.516$ \\
    $1.00$ & $0.500$ & $0.500$ \\
    \bottomrule
  \end{tabular}
\end{table}

The monotone descent in \Cref{tab:redundancy} is not an accident; it holds in
general.

\begin{proposition}[Diversification monotonicity]
\label{prop:monotone}
Let $T$ be a subgroup of equal-ability assets, independent of the rest of the
field, with common pairwise correlation $\rho$ within $T$. Then the subgroup's
combined weight $\sum_{i\in T} w_i$ is non-increasing in $\rho$, falling from its
independent value at $\rho = 0$ to a single asset's weight at $\rho = 1$.
\end{proposition}
\begin{proof}
Write the within-group performances in common-factor form
$X_i = a + \sqrt{\rho}\,Z + \sqrt{1-\rho}\,\varepsilon_i$ ($i \in T$), with
$Z, \varepsilon_i$ independent standard normals. By Slepian's inequality
\parencite{slepian1962}, raising the off-diagonal correlations of a Gaussian
vector while fixing its marginals raises $\Prob(\min_{i\in T} X_i > t)$ for every
$t$, so $\min_{i\in T} X_i$ is stochastically increasing in $\rho$. Since $T$ is
independent of the remaining field, the subgroup wins exactly when its minimum
falls below the ($\rho$-independent) minimum of the complement; that probability
is therefore non-increasing in $\rho$. At $\rho = 0$ the group has $|T|$
independent chances to win; at $\rho = 1$ it collapses to one
(\Cref{prop:redundancy} is the case $|T| = 2$, $\rho = 1$).
\end{proof}

\noindent This is the diversification content of the tilt: correlated clusters
are progressively de-weighted as their internal correlation rises, with no
explicit risk term.

\subsection{The implied objective}
\label{sec:objective}

The tilt is defined procedurally, by winning probabilities. It is natural to ask
what objective, if any, it optimizes. It optimizes a recognizable one.

Write $\theta = -a$ (so larger $\theta$ denotes a stronger competitor, since the
minimum wins) and define the \emph{expected-best} potential
\begin{equation}
  G_C(\theta) = \E\Big[\max_i\,(\theta_i + \eta_i)\Big], \qquad
  \eta \sim \mathcal{N}(0, C).
  \label{eq:surplus}
\end{equation}

\begin{theorem}[Implied objective]
\label{thm:objective}
$G_C$ is convex and $C^\infty$ in $\theta$, the weights are its gradient,
$w = \nabla G_C(\theta)$, and equivalently $w$ is the unique solution of the
regularized program
\begin{equation}
  w \;=\; \argmaxop_{p \in \Delta}\ \big\{\, \langle \theta, p \rangle
  - \Omega_C(p) \,\big\},
  \label{eq:regularized}
\end{equation}
where $\Omega_C = G_C^{*}$ is a convex regularizer (the Fenchel conjugate of
$G_C$) supported on the simplex $\Delta$.
\end{theorem}
\begin{proof}
$G_C$ is an expectation of functions affine in $\theta$, hence convex; with an
almost-surely unique maximizer the envelope theorem gives
$\partial G_C/\partial\theta_i = \Prob(i = \argmaxop_k(\theta_k + \eta_k)) = w_i$,
i.e.\ $w = \nabla G_C$ (the Williams--Daly--Zachary identity
\parencite{williams1977, mcfadden1981}). Equation~\eqref{eq:regularized} is the
Fenchel--Young duality for the convex $G_C$.
\end{proof}

So the portfolio maximizes an ability-weighted exposure $\langle\theta,p\rangle$
minus a convex \emph{diversification penalty} $\Omega_C$ that is determined by
the correlation---a mean--(convex-risk) allocator whose feasible set is the
simplex, which is why the solution is long-only and fully invested. This is the
random-utility ``social-surplus'' / perturbed-optimizer construction
\parencite{mcfadden1981, williams1977, berthet2020}. The choice of perturbation
is decisive: Gumbel noise yields $\Omega = $ Shannon entropy and $w = $ softmax,
a \emph{correlation-blind} (Lucian) allocator; Gaussian noise with covariance
$C$ makes $\Omega_C$ \emph{correlation-aware}, which is precisely what produces
the redundancy consistency of \Cref{prop:redundancy}. Finally, since
$w = \nabla G_C(\theta)$ is the gradient of a convex potential, it is
automatically smooth---consistent with \Cref{thm:smooth}. (For general $C$ the
regularizer $\Omega_C$ has no closed form, unlike the entropy/softmax special
case, but it provably exists and is convex, which is all we use.)

\section{Smoothness and Turnover}
\label{sec:smoothness}

In the online (rebalancing) setting the correlation estimate moves slowly from
one period to the next, and we want the portfolio to move slowly with it:
turnover should track \emph{genuine} change in the correlation, not estimation
or sampling noise. We show the weight map is smooth, with a turnover bound, and
that a common-seed Monte-Carlo realisation inherits it.

Fix the abilities $a$ (calibrated once; they drift slowly) and view the weights
as a function of the correlation matrix $C$. Centring performances by $a$, asset
$i$ wins iff $Y = X - a \sim \mathcal{N}(0, C)$ lands in the \emph{fixed}
polyhedral cone
\begin{equation}
  A_i = \{\, y : y_i - y_j \le a_j - a_i \ \ \forall j \,\},
  \label{eq:wincone}
\end{equation}
so the weight is
$w_i(C) = \Prob_{Y\sim\mathcal{N}(0,C)}(Y\in A_i) = \int_{A_i}\varphi_C(y)\,dy$,
with $\varphi_C$ the $\mathcal{N}(0,C)$ density.

\begin{theorem}[Smoothness]
\label{thm:smooth}
On the open cone of positive-definite correlation matrices, $w_i(C)$ is
$C^\infty$ in $C$ (and in $a$). Moreover, on any compact subset whose smallest
eigenvalue is bounded below by $\lambda > 0$, the map is Lipschitz:
\begin{equation}
  \|w(C_1) - w(C_0)\|_1 \;\le\; L_\lambda \,\|C_1 - C_0\|,
  \qquad L_\lambda = O(1/\lambda).
  \label{eq:lip}
\end{equation}
\end{theorem}

\begin{proof}
The region $A_i$ in \eqref{eq:wincone} is independent of $C$. On the
positive-definite cone $\det C$ and $C^{-1}$ are real-analytic, so
$(C, y) \mapsto \varphi_C(y)$ is $C^\infty$ in $C$; its $C$-derivatives are
polynomials in $y$ times $\varphi_C(y)$, which are integrable and locally
dominated uniformly on compact sets. Differentiation under the integral sign
then gives $w_i \in C^\infty$. For the gradient, Price's theorem
\parencite{price1958}---the Gaussian identity
$\partial \varphi_C/\partial C_{jk} = \partial^2 \varphi_C/\partial y_j\,\partial y_k$
for $j \neq k$---followed by the divergence theorem over the cone $A_i$ rewrites
$\partial w_i/\partial C_{jk}$ as an integral of $\varphi_C$ over the boundary
faces of $A_i$ (the pairwise-tie hyperplanes), which is finite. The gradient is
therefore bounded; since the conditioning of $C$ controls the tie-boundary
density, the bound is $O(1/\lambda)$, giving \eqref{eq:lip}.
\end{proof}

The Lipschitz constant degrades as $C$ approaches singularity
($\lambda \to 0$)---the expected caveat that near-degenerate correlation makes
ties, and hence weights, sensitive. This motivates keeping $C$ well-conditioned
and regenerating the path ensemble when its effective sample size collapses.

\paragraph{Common-seed transport.}
We realise $w(C)$ by Monte Carlo with a \emph{fixed} set of standard-normal
seeds $Z_1, \dots, Z_M$ and the symmetric square root $S(C) = C^{1/2}$ (itself
$C^\infty$ on the positive-definite cone): the paths are
$\widehat{Y}_m(C) = S(C)\,Z_m$ and
$\widehat{w}_i(C) = M^{-1}\sum_m \mathbf{1}[\widehat{Y}_m(C) \in A_i]$. Each
indicator is piecewise constant in $C$, flipping only when a path crosses a face
of $A_i$, so a step $\Delta C$ flips $O(M\,\|\Delta C\|)$ paths and
$\|\Delta\widehat{w}\|_1 = O(\|\Delta C\|)$. Crucially, when $\Delta C = 0$ the
same seeds give the same winners and $\Delta\widehat{w} = 0$ \emph{exactly},
whereas redrawing fresh samples each period incurs an irreducible
$O(1/\sqrt{M})$ turnover whether or not $C$ changed. The common seeds make the
sampling error a fixed offset that cancels in differences; this is the mechanism
behind the incremental (online) rebalancing update.

\Cref{tab:turnover} confirms both statements numerically along a $60$-step
convex path between two random correlation matrices.

\begin{table}[h]
  \centering
  \caption{Turnover along a $60$-step correlation path ($n = 8$,
  $M = 4\times10^4$). Common-seed transport tracks genuine correlation change;
  fresh resampling carries a noise floor that does not vanish when the
  correlation is held static.}
  \label{tab:turnover}
  \begin{tabular}{lc}
    \toprule
    Quantity & Value \\
    \midrule
    mean $\|\Delta C\|_1$ per step                          & $0.372$ \\
    mean $\|\Delta w\|_1$, common-seed transport            & $0.0043$ \\
    mean $\|\Delta w\|_1$, fresh resampling                 & $0.0145$ \\
    Lipschitz ratio $\|\Delta w\|_1 / \|\Delta C\|_1$ (max) & $0.027$ \\
    static $\Delta C = 0$: transport                        & $0$ (exact) \\
    static $\Delta C = 0$: resampling                       & $0.0157$ \\
    \bottomrule
  \end{tabular}
\end{table}

\section{The Indispensable Markets Hypothesis}
\label{sec:capm}

\paragraph{Why capitalization weighting is supposed to be optimal.}
The contradiction lives in the premises, so we recall the argument. Under the
CAPM every investor is a mean--variance optimizer who agrees on the expected
returns $\mu$ and---the load-bearing assumption---on the \emph{same} covariance
$\Sigma$. Two-fund separation then has every investor hold the risk-free asset
together with the \emph{same} tangency portfolio, proportional to
$\Sigma^{-1}(\mu - r\mathbf{1})$. Because all investors hold the identical risky
portfolio, market clearing forces it to equal the supply---the
capitalization-weighted market---so $\mathrm{cap} \propto \Sigma^{-1}(\mu -
r\mathbf{1})$ and the market portfolio is efficient. Capitalization weighting is
optimal \emph{because the covariance is shared}: it is the single portfolio on
which all investors, reasoning from one $\Sigma$, agree.

\paragraph{Why it fails.}
Two premises do not survive contact with practice. First, $\Sigma$ is not known
but estimated, and usable estimates must be \emph{shrunk} toward a structured
target \parencite{ledoitwolf2004}; investors using different, shrunken
covariances hold different tangency portfolios whose aggregate is generically
\emph{not} the capitalization-weighted market. The very step that makes covariance
estimation work destroys the homogeneity the optimality rests on. Second, even
with shared beliefs the clearing argument closes only over the \emph{entire} asset
world: the cap weights inside any investable subset are inherited from the
full-market equilibrium rather than solved for within the subset, and are
generically inefficient there. The rest of this section makes that second failure
precise and structural---it is the Lucian nature of capitalization weighting.

\paragraph{Capitalization weighting is Lucian.}
Write $w_i(S)$ for the weight a rule assigns to asset $i$ within an investable
universe $S$. Capitalization weighting sets
$w_i(S) = \mathrm{cap}_i / \sum_{j\in S}\mathrm{cap}_j$, so the ratio
$w_i(S)/w_j(S) = \mathrm{cap}_i/\mathrm{cap}_j$ is \emph{independent of $S$}. This
is exactly Luce's choice axiom \parencite{luce1959}: the relative weight of two
assets is unaffected by what else is in the universe (independence of irrelevant
alternatives). We call such a rule \emph{Lucian}.

\paragraph{Efficiency is not.}
Mean--variance efficiency lacks this property. The tangency weights on $S$ are
$w^{*}(S) \propto \Sigma_S^{-1}(\mu_S - r\mathbf{1})$, whose ratio
$w^{*}_i(S)/w^{*}_j(S)$ depends on the entire inverse covariance $\Sigma_S^{-1}$
and therefore changes when an asset is added to or removed from $S$ (unless
$\Sigma_S$ is diagonal). Efficiency is non-Lucian.

\begin{proposition}[Lucian weighting is inconsistent under restriction]
\label{prop:restriction}
Capitalization weighting is the Lucian rule with values $\mathrm{cap}_i$. If
returns are correlated (non-diagonal $\Sigma$), then generically a Lucian rule
coincides with mean--variance efficiency on at most one universe in a nested
family. Under the CAPM that universe is the full asset world $U$; on any strict
sub-universe $S \subsetneq U$ the capitalization-weighted portfolio is generically
inefficient, and the discrepancy is governed by the heterogeneity of the assets'
covariances with the excluded set $U \setminus S$.
\end{proposition}
\begin{proof}[Proof sketch]
The cap-weight ratio $\mathrm{cap}_i/\mathrm{cap}_j$ is constant in $S$, while the
efficient ratio $[\Sigma_S^{-1}(\mu_S - r\mathbf{1})]_i / [\,\cdot\,]_j$ varies
with $S$ through $\Sigma_S^{-1}$; a constant and a (generically) non-constant
function of $S$ agree on at most one $S$. CAPM equilibrium fixes equality at $U$.
The wedge on $S$ is the part of each $\mathrm{cap}_i$ reflecting asset $i$'s
covariance with $U \setminus S$---its beta to the excluded assets---which is
irrelevant to the $S$-frontier yet baked into its full-market weight.
\end{proof}

This is the choice-theoretic form of Roll's critique \parencite{roll1977}: the
only testable content of the CAPM is the efficiency of the \emph{true} market
portfolio, and no proxy on a restricted universe inherits it
\parencite{prono2009}. The contribution here is the lens---capitalization
weighting is \emph{Lucian}, efficiency is not, and the contradiction is forced
the moment the universe is restricted. \Cref{prop:redundancy} is the same failure
seen from inside the set: a perfectly correlated duplicate is the extreme case of
``correlation with the rest,'' which a Lucian rule over-weights. (The
fundamental-indexing literature reaches cap-weighting's sub-optimality by a
different route, price inefficiency \parencite{hsu2006}.)

\paragraph{The hypothesis.}
We state the \emph{Indispensable Markets Hypothesis} as the claim that within any
restricted, investable universe the correlation-aware Thurstonian race recovers a
weighting that capitalization---being Lucian---cannot, and that tilting toward it
earns a higher long-run risk-adjusted return than the cap-weighted benchmark. By
\Cref{prop:restriction} the claim is falsifiable: the advantage should grow with
the heterogeneity of the assets' correlations (within $S$ and with the world
outside it) and vanish as that structure flattens. \Cref{sec:experiments}
describes the tests.

\section{Empirical Protocol and Studies}
\label{sec:experiments}

\paragraph{Protocol.}
For each rebalancing date we take capitalization weights as the target
$w^{\mathrm{target}}$, estimate the return correlation $\widehat{C}$ from a
trailing window, calibrate abilities under the reference $C_{\mathrm{calib}}$
(diagonal or one-factor market), and tilt under
$C_{\mathrm{tilt}} = (1-\phi)\,C_{\mathrm{calib}} + \phi\,\widehat{C}$ over a grid
of $\phi$, transporting a fixed seed ensemble between dates. Weights are held
until the next rebalance. We report annualized return, an adjusted Sharpe ratio,
and maximum drawdown, benchmarked against equal-weight and capitalization-weight
portfolios.

\paragraph{Study 1: REITs.}
A panel of real-estate investment trusts (daily returns and market caps) is used
to compare the ability tilt against equal- and capitalization-weighted
baselines across sample budgets $M\in\{10^3,10^5\}$ and correlation scales
$\phi\in\{0,\tfrac12,1\}$. \emph{Results to be inserted.}

\paragraph{Study 2: Russell indices with AI-benefit tilts.}
Russell 1000/2000/3000 constituents are scored by language models for expected
benefit from AI adoption; the scores define the field that is then allocated by
the ability tilt and compared with capitalization weighting and with simple
top-quantile score tilts. \emph{Results to be inserted.}

\begin{table}[h]
  \centering
  \caption{Backtest summary (placeholder).}
  \label{tab:results}
  \begin{tabular}{lccc}
    \toprule
    Portfolio & Ann.\ return & Adj.\ Sharpe & Max drawdown \\
    \midrule
    Equal weight                          & --- & --- & --- \\
    Cap weight (target, $\phi{=}0$)        & --- & --- & --- \\
    Ability tilt ($\phi{=}\tfrac12$)       & --- & --- & --- \\
    Ability tilt ($\phi{=}1$)              & --- & --- & --- \\
    \bottomrule
  \end{tabular}
\end{table}

\section{Related Work}
\label{sec:related}

\paragraph{Choice and ranking.}
The Thurstonian (probit) model \parencite{thurstone1927} and the Luce
\parencite{luce1959} / Bradley--Terry \parencite{bradleyterry1952} /
Plackett--Luce \parencite{plackett1975} family are the two classical accounts of
probabilistic choice. They differ precisely on independence of irrelevant
alternatives, which the Luce family imposes and the Thurstonian model---through
correlated utilities---does not. Our allocator is the Thurstonian side of that
dichotomy applied to portfolio weights.

\paragraph{The racing inverse.}
Assigning winning probabilities to multi-entrant contests and inverting them to
relative abilities is treated by \textcite{harville1973} and, in the lattice form
we use, by \textcite{cotton2021}; the latter supplies the linear-time
forward/inverse maps and tie handling that drive calibration.

\paragraph{Capitalization weighting and its critics.}
The CAPM \parencite{sharpe1964} underwrites capitalization weighting; Roll's
critique \parencite{roll1977} and its correlation-aware refinements
\parencite{prono2009} show that the justification requires the entire asset
universe---which we recast as a clash between the Lucian invariance of cap
weighting and the non-Lucian nature of efficiency. Fundamental indexing
\parencite{hsu2006} reaches cap weighting's sub-optimality from price
inefficiency, while shrinkage \parencite{ledoitwolf2004} and tracking-error
analysis \parencite{roll1992} document the fragility of the estimated inputs. The
classical optimization baseline is \textcite{markowitz1952}.

\paragraph{Perturbed optimizers.}
The identity ``weights equal the gradient of an expected maximum'' is the
random-utility social-surplus construction \parencite{mcfadden1981, williams1977}
and, in machine learning, the perturbed-optimizer / differentiable-argmax
framework \parencite{berthet2020}. Gumbel perturbations give the
entropy-regularized softmax; our Gaussian, correlation-bearing perturbation gives
the correlation-aware regularizer of \Cref{thm:objective}.

\section{Discussion}
\label{sec:discussion}

The ability tilt is a small idea with useful properties: it is long-only and
feasible without optimization, it never inverts a covariance, it reproduces a
chosen benchmark by construction and tilts only in response to correlation, and
it solves an interpretable correlation-aware regularized program
(\Cref{thm:objective}). The smoothness of the weight map (\Cref{thm:smooth})
bounds turnover and licenses both cheap online updates and gradient-based tuning.
Open questions include the choice of reference correlation $C_{\mathrm{calib}}$
(diagonal versus market) and of $\phi$; the treatment of estimation error in the
abilities (a Thurstonian analogue of covariance shrinkage); and richer
perturbations---a temperature controlling concentration, fat-tailed
performances, and using $C_{\mathrm{tilt}}$ to express an explicit correlation
view rather than the sample estimate.

\section{Conclusion}
\label{sec:conclusion}

Thurstone portfolios weight assets by their probability of winning a correlated
race, with abilities calibrated so that the race reproduces a chosen benchmark at
a reference correlation. The portfolio is then a controlled perturbation of the
benchmark, driven only by how the actual correlation departs from the reference.
The construction is long-only and inversion-free; it reproduces the benchmark
exactly when the two correlations agree; it is redundancy-coherent where
capitalization weighting, being Lucian, is not; it solves a correlation-aware
regularized program; and its weights vary smoothly with the correlation, which
bounds turnover and supports online rebalancing at scale. Whether the tilt pays
for itself out of sample is the empirical question the protocol and studies above
are designed to answer.

\printbibliography

\end{document}
